%% 中文摘要页
\pdfbookmark[0]{摘要}{abstract-zh}
\begin{center}\heiti\sanhao\textbf{
		摘要
	}
\end{center}

为了适应数据要素化背景下数据流通需求场景，
针对分布式环境下多节点参与，数据安全授权变得尤为重要，
本项目提出分布式环境下基于国密的代理重加密算法，
将传统代理重加密中的单代理节点，拆分为多个代理节点，
结合门限代理重加密和国密算法的特点，确保在客户端与中心服务器间、中心服务器与不同节点间、
客户端和不同节点间的数据传输和存储过程中，数据保持完整和隐私性。

与传统的访问控制和权限管理方法相比，
本项目算法具有以下优点：
高强度加密：使用门限代理重加密技术对数据进行高强度加密，保护数据的机密性和隐私性。
分布式架构：采用分布式架构实现对分布式数据的访问控制和授权管理，解决了传统方法在分布式场景下的难题。
自主可控：基于国密算法设计，保证了数据的自主可控性，可以有效地避免数据泄露和信息安全风险。
高效性能：算法实现简单、计算量小，可在现有的计算资源下高效地进行数据安全共享和授权管理。

为了验证所提出系统的有效性和性能，本项目进行了一系列的实验测试。
测试结果表明，该基于国密算法的代理重加密系统在分布式环境中具有优越的加密速度、解密效率和扩展性。

因此，本项目方案为分布式数据的安全传输和存储提供了一个高效且安全的解决方案，
并为区块链环境下隐私计算带来了新的研究视角和实践经验。

\textbf{关键词：}分布式计算；国密算法；代理重加密；数据共享；权限管理
\clearpage
